\documentclass[12pt,a4paper]{ctexart}
\usepackage[top=2.5cm,bottom=2.5cm,left=3cm,right=3cm]{geometry}
\usepackage{graphicx}
\usepackage{fancyhdr}
\usepackage{titlesec}
\usepackage{enumitem}
\usepackage{xcolor}
\usepackage{booktabs}
\usepackage{listings}
\usepackage{hyperref}

% 去除默认目录名
\ctexset{contentsname={}}

% 页眉页脚
\pagestyle{fancy}
\fancyhf{}
\fancyhead[L]{\small 审计智鉴流式数据处理系统 V1.0}
\fancyhead[R]{\small \leftmark}
\fancyfoot[C]{\thepage}
\renewcommand{\headrulewidth}{0.4pt}
\renewcommand{\footrulewidth}{0pt}
\renewcommand{\sectionmark}[1]{\markboth{\thesection\quad #1}{}}

% 章节标题样式
\titleformat{\section}{\heiti\zihao{-2}\bfseries}{\thesection}{1em}{}
\titleformat{\subsection}{\heiti\zihao{3}\bfseries}{\thesubsection}{1em}{}
\titleformat{\subsubsection}{\heiti\zihao{4}\bfseries}{\thesubsubsection}{1em}{}

% 列表样式
\setlist[itemize]{leftmargin=2em,itemsep=0.3em}
\setlist[enumerate]{leftmargin=2em,itemsep=0.3em}

% 代码样式
\lstset{
    language=Python,
    basicstyle=\ttfamily\small,
    frame=single,
    breaklines=true,
    showstringspaces=false,
    tabsize=4
}

% 自定义颜色
\definecolor{hintblue}{RGB}{37,99,235}
\definecolor{warnorange}{RGB}{245,158,11}
\definecolor{successgreen}{RGB}{16,185,129}

% 截图占位框
\newcommand{\screenshotplaceholder}[2]{%
    \begin{center}
    \fbox{\parbox{0.85\textwidth}{%
        \centering
        \vspace{1.5cm}
        {\large \textbf{【截图位置：#1】}}\\[0.5em]
        {\small #2}
        \vspace{1.5cm}
    }}
    \end{center}
}

\begin{document}

% ==================== 封面 ====================
\begin{titlepage}
\centering
\vspace*{3cm}
{\Huge\heiti\bfseries 审计智鉴流式数据处理系统}\\[0.8cm]
{\LARGE\heiti 操作说明书}\\[0.5cm]
{\Large V1.0}\\[3cm]
\begin{tabular}{rl}
\textbf{软件全称：} & 审计智鉴流式数据处理系统 V1.0 \\[0.3em]
\textbf{软件简称：} & 审计智鉴系统 \\[0.3em]
\textbf{版本号：} & V1.0 \\[0.3em]
\textbf{编制日期：} & 2026年01月 \\
\end{tabular}
\vfill
\end{titlepage}

% ==================== 目录 ====================
\tableofcontents
\newpage

% ==================== 第1章 系统概述 ====================
\section{系统概述}

\subsection{系统用途与背景}

财务造假一直是资本市场的一大顽疾。不管是早几年的康美药业、康得新，还是后来的瑞幸咖啡，这些案例都在提醒我们：光靠人工去翻几百页的年报，效率实在太低，而且很多风险信号藏在文本细节里，肉眼根本看不出来。

审计智鉴系统这个项目，最初的出发点其实挺简单的——就是把"读财报"这件事交给程序来做。当然不是完全替代审计师，而是提供一个辅助工具，帮审计人员快速定位可疑点，把精力集中在真正需要深挖的地方。

系统核心解决的是两个问题：一是传统财务指标分析（比如存贷双高、现金流背离这些），二是管理层讨论与分析（MD\&A）文本中的语义风险。把这两块结合起来，形成一个相对完整的舞弊风险评估体系。

\subsection{技术架构}

整个系统采用前后端分离的架构，这个选型主要考虑了开发效率和部署灵活性。

后端用的是 \textbf{FastAPI}，这是一个基于 Python 的异步 Web 框架，性能不错，而且天然支持异步 IO，对于调用外部 AI 接口这种场景很合适。后端负责所有核心业务逻辑：用户认证、检测分析、AI 问答、报告生成、会员计费等。

前端用的是 \textbf{Streamlit}，这是一个专门面向数据应用的 Python 框架。说实话，选它而不是 React 或 Vue，主要是因为开发速度快——全是 Python 代码，不用切语言。而且 Streamlit 对数据可视化支持得很好，图表、表格、进度条这些组件都是现成的。当然它也有缺点，就是 UI 自定义能力不如传统前端框架强，不过对于我们这个工具型应用来说已经够用了。

数据层用的是 \textbf{MySQL}，通过 SQLAlchemy ORM 进行数据库操作。缓存用了简单的内存缓存，对于当前的用户量来说暂时够用了。

AI 能力接入的是阿里云百炼（DashScope）的通义千问模型，支持流式输出。之所以选这个，一方面是中文理解能力强，另一方面是 API 接口是 OpenAI 兼容格式，对接起来比较方便。

\subsection{系统核心能力}

概括来说，系统具备以下几项核心能力：

\begin{enumerate}
    \item \textbf{双模融合检测}：传统财务指标占 40\% 权重，AI 文本分析占 60\% 权重，综合计算舞弊概率。
    \item \textbf{七维度文本风险分析}：从语义矛盾、风险披露、语调异常、文本-数据一致性、关联隐藏、信息密度、回避表述七个角度评估 MD\&A 文本。
    \item \textbf{SHAP 可解释性}：基于博弈论的 Shapley 值方法，量化每个特征对最终判断的贡献度，满足审计报告的可解释性要求。
    \item \textbf{流式 AI 问答}：采用 SSE（Server-Sent Events）技术，AI 回答内容逐字实时呈现，而不是等全部生成完再一次性显示。
    \item \textbf{多格式报告导出}：检测结果支持导出为 PDF、Word、Excel 三种格式，方便存档和汇报。
\end{enumerate}

\newpage

% ==================== 第2章 运行环境 ====================
\section{运行环境说明}

\subsection{服务器环境}

本系统部署于云服务器环境，操作系统采用 Linux Ubuntu 22.04 LTS。之所以选这个版本，主要是因为稳定性好，社区支持也成熟，出问题的时候网上能找到解决方案。

服务器配置要求并不高，前期 2 核 CPU、4GB 内存就能跑起来。如果用户量上来了，主要瓶颈在外部 AI 接口的调用并发上，而不是系统本身的计算压力。毕竟核心的 XGBoost 模型推理很快，毫秒级就能完成。

\subsection{访问方式}

系统无专属域名，用户通过公网 IP 地址加端口号的方式访问使用。后端 API 服务运行在 8000 端口，前端 Streamlit 界面运行在 8501 端口。这种部署方式省去了域名注册和备案的流程，对于一个小型工具类应用来说比较省事。

前后端分离架构下，前端页面加载后，所有数据交互都通过 RESTful API 完成。用户在前端上传文件、提交检测请求，前端把请求发到后端，后端处理完再返回结果，前端负责渲染展示。

\subsection{运行依赖}

系统运行需要以下环境：

\begin{itemize}
    \item Python 3.9 或更高版本
    \item MySQL 8.0 数据库
    \item 阿里云 DashScope API 访问权限（用于 AI 分析）
    \item 足够的磁盘空间用于存储上传文件和生成的报告
\end{itemize}

数据库初始化在系统启动时自动完成，会创建所有必要的表结构。第一次启动时可能需要等几秒，因为 SQLAlchemy 在检查和同步表结构。

\newpage

% ==================== 第3章 系统功能模块 ====================
\section{系统功能模块}

本章对系统的五个主要功能模块进行说明。每个模块都围绕"数据接收—运算处理—结果输出"的流程展开，前后端通过接口联动完成整个业务闭环。

\subsection{用户认证模块}

这个模块负责用户的注册、登录和权限管理。密码采用 PBKDF2-SHA256 算法哈希存储，不是明文也不是简单的 MD5，安全性上有基本保障。登录成功后后端签发 JWT Token，有效期 7 天，前端把这个 Token 存在 Cookie 里，后续请求自动带上。

用户分为免费版、专业版、企业版三个等级，不同等级在检测次数和 AI 问答次数上有差异。免费版每月 3 次检测额度，每天 5 次 AI 问答；专业版检测不限次数，AI 问答每天 50 次；企业版全部不限量。这个分级主要是为了覆盖不同用户的需求，个人审计师用免费版体验一下，机构用户买专业版或企业版。

\subsection{舞弊检测模块}

这是系统的核心模块。用户可以通过三种方式提交检测：

\begin{enumerate}
    \item \textbf{文件上传}：上传 Excel 或 CSV 格式的财务数据文件，系统会自动解析年份和各项指标。还可以补充上传 MD\&A 文本文件（txt、docx、pdf 都可以），用于文本风险分析。
    \item \textbf{内置案例库}：系统预置了若干经典案例，比如康美药业、贵州茅台等，用户可以直接加载这些案例数据进行检测，方便快速体验功能。
    \item \textbf{手动录入}：在页面上直接填写各项财务指标，适合数据量不大的场景。
\end{enumerate}

检测流程大概是这样的：后端收到请求后，先进行传统财务指标分析（存贷双高、现金流背离、存货异常等），然后调用 AI 模型对 MD\&A 文本进行七维度分析，最后把两部分得分加权汇总，计算出一个 0 到 100 的风险评分和舞弊概率。整个过程大概 3 到 5 秒，取决于外部 AI 接口的响应速度。

\subsection{AI 问答模块}

这是一个流式输出的智能问答功能。用户在输入框里提问，比如"什么是存贷双高"或者"帮我分析一下康美药业的案例"，系统会通过 SSE 流式返回 AI 的回答，文字是逐字出现的，有一种"正在打字"的实时感。

流式输出在技术实现上用了 Server-Sent Events，前端建立一个长连接，后端通过 httpx 的 stream 模式读取 AI 接口的流式响应，然后一段一段地推送给前端。这比传统的"等全部内容生成完再返回"体验要好很多，用户不用干等着。

系统还预设了五个类别的问题推荐：理论类、实践类、政策类、案例类、平台使用类。侧面栏里可以直接点这些问题，省得打字。

\subsection{财务助手模块}

财务助手下面分了三个子功能：

\begin{itemize}
    \item \textbf{税务测算}：按最新个税法和企业所得税法计算应纳税额。个税测算支持月薪、年终奖、五险一金、专项附加扣除等，还能对比年终奖"单独计税"和"合并计税"两种方式哪种更省税。企业税测算支持小规模纳税人和一般纳税人两种类型，不同行业的增值税税率也做了区分。
    \item \textbf{简易报税辅助}：用户可以上传增值税发票图片，系统调用多模态 AI 模型（qwen-vl-max）识别发票内容，提取金额、税率、税额等信息，然后自动生成报税数据汇总表。这个功能目前主要支持增值税发票，其他类型单据的识别准确率会有波动。
    \item \textbf{财税咨询}：和 AI 问答类似，但侧重点在税务政策解读和实操建议上。同样是流式输出。
\end{itemize}

\subsection{报告管理模块}

检测完成后，用户可以生成正式的检测报告。报告类型分基础版、专业版、企业版三种，内容深度不同。基础版主要是检测结果摘要和风险标签；专业版会加上 SHAP 特征分析解读、整改建议、IPO 对标分析；企业版更详细，包含完整的技术细节和原始数据附录。

报告支持导出为 PDF、Word、Excel 三种格式。PDF 用的是预定义好的模板，排版比较正式；Word 方便用户二次编辑；Excel 则是把数据表格单独导出来，方便做进一步分析。

\newpage

% ==================== 第4章 系统操作流程 ====================
\section{系统操作流程}

本章通过实际操作流程，说明系统的使用方法。以下所有操作均在前端交互界面完成，后端通过 API 接口实时响应。

\subsection{首次使用流程}

第一次打开系统，看到的是首页。页面上方有一条渐变色装饰条，中间是品牌 Logo 和导航菜单，整体布局比较简洁。首页主体是一个深色背景的大图区，上面有系统的核心定位说明，下面是一排动态统计数字（AI 识别准确率、已分析财报数量等），这些数字在页面滚动到可视区域时会有从 0 滚动到目标值的动画效果。

\screenshotplaceholder{img-01}{此处应插入系统主界面截图，需包含顶部导航栏、Hero 区域和动态统计数据栏}

新用户需要先注册账号。点击右上角的"登录 / 注册"按钮，弹出登录弹窗。弹窗分登录和注册两个标签页，注册时填写用户名、密码、邮箱（可选）、手机号（可选）即可。注册成功后自动登录，跳转到首页。

\screenshotplaceholder{img-02}{此处应插入登录/注册弹窗截图}

\subsection{执行舞弊检测}

登录后，点击导航栏的"舞弊检测"进入检测页面。页面上方有企业名称和证券代码的输入框，下面是年份区间设置，再下面是文件上传区域。

\screenshotplaceholder{img-03}{此处应插入舞弊检测-文件上传界面截图}

以文件上传方式为例，操作步骤如下：

\begin{enumerate}
    \item 填写企业名称（必填）和证券代码（可选）。
    \item 设置数据年份区间，系统默认 2020 到 2023 年。
    \item 上传财务数据文件，支持 xlsx、xls、csv 格式。如果文件里包含"年份"列，系统会自动按年份分组解析；如果没有年份列，则把整个文件当作起始年份的数据。
    \item （可选）上传 MD\&A 文本文件，支持 txt、docx、doc、pdf 格式，可以多选。
    \item 点击"解析文件内容"按钮，系统会解析文件并在页面上显示预览结果，包括各年度提取的财务指标。
    \item 确认数据无误后，点击"开始批量检测"按钮，系统会逐年度执行检测分析，页面上会显示进度条。
\end{enumerate}

检测完成后，页面自动跳转到结果展示页。结果页顶部有一个综合信息卡片，显示舞弊概率、风险评分、风险标签数量、置信度四个核心指标。下面左侧是舞弊概率的仪表盘图表（环形图），右侧是风险标签统计。

\screenshotplaceholder{img-04}{此处应插入舞弊检测-检测结果截图，需包含风险仪表盘和风险标签}

再往下是 AI 风险特征雷达图，七个维度（语义矛盾度、风险披露完整性、异常乐观语调、文本-数据一致性、关联隐藏指数、信息密度异常、回避表述强度）的得分一目了然。右侧是 SHAP 特征重要性排行，用进度条的形式展示 Top 5 特征对模型判断的贡献度。

\subsection{使用 AI 问答（流式输出效果）}

点击导航栏的"AI 问答"进入问答页面。如果是第一次进入，页面中央会显示一个欢迎卡片，介绍这是一个财务舞弊领域的专业问答助手。

在底部的输入框里输入问题，比如"康美药业舞弊案的关键识别点是什么"，按回车发送。这时页面会显示用户的问题气泡，然后 AI 助手的位置会出现三个跳动的小圆点（ thinking 动画），表示后端正在处理。

大概一两秒后，回答文字开始逐字出现——这就是 SSE 流式输出的效果。文字不是一下子全部跳出来，而是像有人在实时打字一样，一个一个字符地显示出来。这种体验和传统网页的"提交表单→等待加载→整页刷新"完全不同，更接近一个正在运行的应用程序。

\screenshotplaceholder{img-05}{此处应插入 AI 问答流式输出效果截图，需能看到文字逐字出现的状态。这是证明系统不是静态网页的关键截图，必须包含}

流式输出结束后，完整的回答会保留在对话记录中。用户可以接着追问，系统支持多轮上下文对话。所有历史问答记录在侧边栏的"历史记录"里可以查看，也可以收藏常用问答。

\subsection{使用税务测算}

进入"财务助手"页面，点击"税务中心"标签。这里分三个子标签：税务测算、简易报税、财税咨询。

以个税测算为例：填写税前月薪、年终奖金额、每月五险一金缴纳额、专项附加扣除额，然后点击"计算个税"按钮。系统会立即计算出年收入、年扣除额、适用税率、年个税、税后年收入五个指标。下面还有一个展开区域，显示应纳税所得额的详细计算过程，包括基本减除费用 6 万元、五险一金年总额、专项附加扣除年总额等逐项明细。

如果填写了年终奖，系统还会自动对比"单独计税"和"合并计税"两种方式，推荐更省税的方案，并显示能节省多少税款。这个功能对于年底拿年终奖的人来说挺实用的。

\screenshotplaceholder{img-06}{此处应插入税务测算结果截图}

\subsection{生成并导出报告}

检测完成后，在"我的检测"页面可以看到所有历史检测记录。点击某条记录旁边的"查看详情"，进入详细结果页。页面上方有一个"导出报告"按钮，点击后选择报告类型（基础版/专业版/企业版）和导出格式（PDF/Word/Excel），系统会在后端生成文件并提供下载链接。

\screenshotplaceholder{img-07}{此处应插入报告导出界面截图}

报告生成大概需要 5 到 10 秒，因为要把所有检测数据、图表、分析文本整合到一个文档里。PDF 版本用的是预定义模板，带有页眉页脚和封面，比较正式，可以直接用于工作汇报。

\subsection{会员订阅}

免费版用户的检测额度用完后，系统会提示升级会员。点击导航栏的"会员"进入定价页面。页面展示了三个版本的会员方案：基础版（299元/年）、专业版（999元/年）、企业版（按需定价）。每个方案下面列出了包含的功能和额度。

\screenshotplaceholder{img-08}{此处应插入会员中心定价页面截图}

用户选择一个方案后，点击"立即订阅"按钮，系统创建订单并跳转到支付页面。目前支持支付宝和微信支付两种渠道（支付功能为模拟演示，实际部署时可接入真实支付接口）。支付成功后会员权益立即生效，无需重新登录。

\newpage

% ==================== 第5章 核心代码逻辑说明 ====================
\section{核心代码逻辑说明}

\subsection{前后端联动流程}

整个系统的前后端交互遵循标准的请求-响应模式，但 AI 问答模块采用了长连接的流式传输，这里单独说明一下。

普通接口（如用户登录、获取检测历史）的流程很常规：前端用 Python 的 requests 库发送 HTTP 请求，后端 FastAPI 处理请求、操作数据库、返回 JSON 数据。前后端之间的数据格式统一用 Pydantic Schema 进行校验，避免类型错误。

AI 问答的流式接口则不同。前端发送 POST 请求时，虽然也是用 requests，但加上了 \texttt{stream=True} 参数。后端收到请求后，不是一次性生成完整回答再返回，而是调用 httpx 的异步流式客户端，向阿里云 DashScope API 发起流式请求。DashScope 返回的响应是一段一段的 SSE 数据（格式为 \texttt{data: \{...\}}），后端把这些数据片段原样转发给前端。前端通过 \texttt{response.iter\_lines()} 逐行读取，解析出其中的文字内容，再用 Streamlit 的 \texttt{st.write\_stream()} 方法实现逐字渲染。

这个流程里有两个关键点：一是 httpx 的异步流式支持，让后端可以一边接收 AI 接口的数据，一边转发给前端，不需要等全部内容；二是 Streamlit 的 \texttt{write\_stream} 方法，它内部做了增量渲染的优化，只更新变化的部分，不会整段文字重绘导致闪烁。

\subsection{风险评分算法}

舞弊概率的计算采用双模融合策略，公式可以简单表示为：

\begin{center}
\fbox{\parbox{0.8\textwidth}{%
\centering
总风险评分 = 基础风险分（35）+ 传统财务风险评分（0--40）+ AI 文本风险评分（0--65）\\[0.3em]
舞弊概率 = 0.50 + (总风险评分 / 100) $\times$ 0.5
}}
\end{center}

传统财务风险评分主要考察三个经典指标：存贷双高（货币资金和短期借款同时处于高位）、现金流与利润背离（净利润为正但经营现金流为负）、存货占比异常（存货占总资产比例过高）。每个指标达到阈值时加上对应的分数，总分封顶 40 分。

AI 文本风险评分则从 MD\&A 文本中提取七个维度的特征得分，每个维度有不同的权重。比如"文本-数据一致性"和"语义矛盾度"的权重最高（2.0 倍），因为这两个维度对舞弊判断的区分度最强。所有维度加权后归一化到 65 分。

基础风险分设为 35 分而不是 0 分，这个设计是有考虑的——财务数据本身就不完全可信，即使没有明显异常，也应该保留一定的基础警惕性。最终舞弊概率被限制在 50\% 到 95\% 之间，最低不会低于 50\%，因为"完全清白"的判断在审计场景下其实是不太负责任的。

\subsection{SHAP 可解释性分析}

模型的可解释性对于审计类应用很重要——你不能只告诉用户"这家企业有 80\% 的舞弊概率"，还得解释清楚为什么是 80\%，哪些因素推动了这个判断。

系统采用 SHAP（SHapley Additive exPlanations）方法来计算每个特征对最终预测结果的贡献度。SHAP 值的核心思想来自博弈论：把每个特征看作一个"参与者"，计算它加入"联盟"时对预测结果的边际贡献。这种方法有一个很好的数学性质，就是所有特征的 SHAP 值加起来等于预测值与基准值的差，满足"可加性解释"。

在实际实现中，由于用的是 XGBoost 模型，可以直接调用 \texttt{shap.TreeExplainer} 来计算 SHAP 值，速度很快。计算结果里既有正值（推动舞弊概率升高）也有负值（推动舞弊概率降低），系统取绝对值最大的前 5 个特征展示给用户，并用进度条的形式可视化。

\newpage

% ==================== 第6章 系统运行总结 ====================
\section{系统运行总结}

审计智鉴系统从立项到现在，功能上已经比较完整了。回头看，有几个地方当初的设计决定现在觉得还是挺对的。

第一个是用 Streamlit 做前端。虽然 UI 自由度不如 React，但开发效率真的高。整个前端就一个 app.py 文件，4900 多行代码，涵盖了所有页面和交互逻辑。如果是传统前后端分离的项目，光前端代码可能就得拆成几十个文件。对于一个小团队来说，这种"单文件搞定一切"的模式反而更实用。

第二个是双模融合的思路。一开始我们也想过完全依赖 AI 做端到端的判断，但测试下来发现纯 AI 方案的稳定性不够，特别是财务数据这种对准确性要求很高的场景，还是需要一些规则兜底。传统指标 + AI 文本分析的组合，既保证了可解释性，又发挥了 AI 在语义理解上的优势。

当然也有一些遗憾。比如报告生成的排版，PDF 版本虽然能看，但和专业排版软件做的还是有差距；再比如票据识别功能，目前主要支持增值税发票，其他类型单据的识别效果还不太稳定。这些算是留给后续版本的改进空间吧。

总的来说，这套系统已经能够独立完成从数据上传、风险分析、结果展示到报告导出的完整流程，部署在云服务器上后可以通过公网 IP 直接访问，不需要额外的域名或备案。对于中小型审计机构或者财务研究人员来说，作为一个辅助分析工具，应该是有一定实用价值的。

\vspace{2cm}
\begin{center}
\rule{0.6\textwidth}{0.4pt}\\[0.5em]
{\small 本文档为审计智鉴流式数据处理系统 V1.0 配套操作说明书}\\
{\small 编制日期：2026年01月}
\end{center}

\end{document}
